% Generated draft: requires mathematical and source review.
\documentclass[11pt]{article}
\usepackage[margin=25mm]{geometry}
\usepackage{fontspec}
\setmainfont{Latin Modern Roman}
\usepackage{amsmath,amssymb,mathtools,unicode-math}
\setmathfont{Latin Modern Math}
\usepackage{enumitem,needspace,xurl,hyperref}
\setlength{\parindent}{0pt}
\setlength{\parskip}{5pt}
\setlength{\emergencystretch}{3em}
\newcommand{\N}{\mathbb{N}}
\newcommand{\Z}{\mathbb{Z}}
\newcommand{\Q}{\mathbb{Q}}
\newcommand{\R}{\mathbb{R}}
\newcommand{\C}{\mathbb{C}}
\newcommand{\F}{\mathbb{F}}
\newcommand{\A}{\mathbb{A}}
\newcommand{\PP}{\mathbb{P}}
\newcommand{\E}{\mathbb{E}}
\newcommand{\eps}{\varepsilon}
\newcommand{\dd}{\,\mathrm{d}}
\begin{document}
\section*{UnsolvedMath ID: 2301015}
\textbf{Source problem code:} catalogue entry
\subsection{Definitions and mathematical statement}
Let $f$ be a nonconstant meromorphic function on $\C$. With $T(r,f)$ its Nevanlinna characteristic and $N(r,a;f)$ the integrated counting function of $a$-points, counted with multiplicity, use
\[
\lambda=\liminf_{r\to\infty}\frac{\log T(r,f)}{\log r},\qquad
\delta(a,f)=1-\limsup_{r\to\infty}\frac{N(r,a;f)}{T(r,f)}.
\]
For $a=\infty$, the counting function counts poles. The source assumes $\delta=\delta(a,f)>0$ and asks for a sequence $r_\nu\to\infty$ along which the angular measure of the part of $|z|=r_\nu$ where $f(z)$ is ``close to $a$'' is at least
\[
B(\lambda,\delta)+o(1),\qquad
B(\lambda,\delta)=\frac{4}{\lambda}\arcsin\sqrt{\frac{\delta}{2}}.
\]
Here angular measure means Lebesgue measure of the corresponding angles in $[0,2\pi)$, not arc length multiplied by $r_\nu$. The real principal inverse sine is intended. The displayed finite expression is defined when $0<\lambda<\infty$; the source does not specify this restriction or conventions at zero or infinite lower order.

To turn ``close'' into a precise assertion, a closeness rule must be fixed. For example, with a prescribed threshold $\eps(r)>0$ and a chosen distance $d$ to $a$, one would define
\[
E(r,a)=\{\theta\in[0,2\pi):d(f(re^{i\theta}),a)<\eps(r)\},
\]
and ask for $\liminf_{\nu\to\infty}|E(r_\nu,a)|\ge B(\lambda,\delta)$. This is an explanatory schema: the source supplies neither the distance nor the threshold or its relation to $T(r,f)$, so no freely chosen or shrinking threshold is silently imposed. A spherical convention would be required to include $a=\infty$ uniformly, while a finite-value Euclidean convention is another possible reading.

There is a further literal qualification: angular measure is at most $2\pi$, whereas the printed expression has no cap and can formally exceed $2\pi$ for unspecified parameter ranges. The extract supplies neither a minimum with $2\pi$ nor a restriction ensuring $B(\lambda,\delta)\le2\pi$. These are unresolved formulation issues; the bound is retained as printed rather than silently repaired.

The supplied summary labels Edrei's spread conjecture solved by Baernstein and gives the references below. That reported resolution must be checked in its precise threshold, lower-order, and saturation conventions before identifying it with the incomplete literal wording here. No proof or independent verification is supplied by this draft.
\subsection{Short English statement}
The spread question asks how large an angular portion of selected large circles has values near a deficient value. The source gives a sharp-looking lower-order bound but omits what ``near'' means and the conventions when the expression exceeds a full circle or the lower order is zero. Its supplied summary attributes the solved spread theorem to Baernstein; the exact theorem needs comparison with those missing conventions.
\subsection{Sources}
\begin{itemize}
\item Local source registry: \url{data/unsolvedmath/problems.json}, permanent UnsolvedMath ID $2301015$, source code catalogue entry, including the supplied solved classification and research summary. Dataset attribution: Ulam AI and UnsolvedMath Contributors; original-source attribution and rights are retained.
\item Hayman, \emph{Research Problems in Function Theory}, source-list date 2018, Problem 1.15, Edrei's spread conjecture: \url{https://arxiv.org/abs/1809.07200}. The metadata reports source-TeX extraction and July 29, 2026 access, with current-status and rights review flags. The original trailing ``see Edrei .'' reference is incomplete in the extract.
\item A. Baernstein II, ``Proof of Edrei's spread conjecture'', Proceedings of the London Mathematical Society (3) 26 (1973), 418--434, as supplied. The summary also cites a Bulletin announcement in Bulletin of the American Mathematical Society 78 (1972), 277--278, without a separate title. No linked source or cited paper was inspected for this draft.
\end{itemize}
\end{document}
